\begin{ex}
Hình tam giác có bao nhiêu cạnh?
\shortans{3}
\loigiai{
Hình tam giác có 3 cạnh.
}
\end{ex}

\begin{ex}
Hình tứ giác có bao nhiêu đỉnh?
\shortans{4}
\loigiai{
Hình tứ giác có 4 đỉnh.
}
\end{ex}

\begin{ex}
Hình chữ nhật có mấy cặp cạnh đối diện dài bằng nhau?
\shortans{2}
\loigiai{
Hình chữ nhật có 2 cặp cạnh đối diện bằng nhau.
}
\end{ex}

\begin{ex}
Hình vuông có bao nhiêu cạnh bằng nhau?
\shortans{4}
\loigiai{
Hình vuông có 4 cạnh dài bằng nhau.
}
\end{ex}

\begin{ex}
Một hình vuông có cạnh dài 7 cm. Tính chu vi hình vuông đó theo đơn vị xăng-ti-mét.
\shortans{28}
\loigiai{
$7 \times 4 = 28$ cm.
}
\end{ex}

\begin{ex}
Một hình chữ nhật có chiều dài 9 cm và chiều rộng 5 cm. Chu vi hình chữ nhật đó là bao nhiêu xăng-ti-mét?
\shortans{28}
\loigiai{
$(9 + 5) \times 2 = 28$ cm.
}
\end{ex}

\begin{ex}
Một hình chữ nhật có chu vi 24 cm. Nửa chu vi hình chữ nhật đó là bao nhiêu xăng-ti-mét?
\shortans{12}
\loigiai{
$24 : 2 = 12$ cm.
}
\end{ex}

\begin{ex}
Một hình chữ nhật có nửa chu vi là 15 cm, chiều rộng là 6 cm. Chiều dài hình chữ nhật đó là bao nhiêu xăng-ti-mét?
\shortans{9}
\loigiai{
$15 - 6 = 9$ cm.
}
\end{ex}

\begin{ex}
Một hình tam giác có độ dài 3 cạnh là 6 cm, 8 cm, 10 cm. Chu vi hình tam giác đó là bao nhiêu xăng-ti-mét?
\shortans{24}
\loigiai{
$6 + 8 + 10 = 24$ cm.
}
\end{ex}

\begin{ex}
Một hình tứ giác có độ dài các cạnh là 5 cm, 7 cm, 8 cm, 10 cm. Chu vi hình tứ giác đó là bao nhiêu xăng-ti-mét?
\shortans{30}
\loigiai{
$5 + 7 + 8 + 10 = 30$ cm.
}
\end{ex}

\begin{ex}
Một hình vuông có chu vi 36 cm. Độ dài một cạnh của hình vuông đó là bao nhiêu xăng-ti-mét?
\shortans{9}
\loigiai{
$36 : 4 = 9$ cm.
}
\end{ex}

\begin{ex}
Một hình vuông có chu vi 20 cm. Độ dài cạnh hình vuông đó là bao nhiêu xăng-ti-mét?
\shortans{5}
\loigiai{
$20 : 4 = 5$ cm.
}
\end{ex}

\begin{ex}
Ghép hai hình vuông cạnh 3 cm lại với nhau tạo thành một hình chữ nhật. Chiều dài hình chữ nhật đó là bao nhiêu xăng-ti-mét?
\shortans{6}
\loigiai{
$3 + 3 = 6$ cm.
}
\end{ex}

\begin{ex}
Ghép hai hình vuông cạnh 4 cm lại thành hình chữ nhật. Chu vi hình chữ nhật tạo thành là bao nhiêu xăng-ti-mét?
\shortans{24}
\loigiai{
Chiều dài = 8 cm, chiều rộng = 4 cm. Chu vi = $(8 + 4) \times 2 = 24$ cm.
}
\end{ex}

\begin{ex}
Một hình chữ nhật có chiều dài 12 cm, chiều rộng bằng $\frac{1}{3}$ chiều dài. Chiều rộng hình chữ nhật đó là bao nhiêu xăng-ti-mét?
\shortans{4}
\loigiai{
$12 : 3 = 4$ cm.
}
\end{ex}

\begin{ex}
Một hình tam giác đều có 3 cạnh bằng nhau, mỗi cạnh dài 9 cm. Chu vi hình tam giác đó là bao nhiêu xăng-ti-mét?
\shortans{27}
\loigiai{
$9 \times 3 = 27$ cm.
}
\end{ex}

\begin{ex}
Một hình tam giác có chu vi 30 cm, biết 3 cạnh bằng nhau. Độ dài mỗi cạnh là bao nhiêu xăng-ti-mét?
\shortans{10}
\loigiai{
$30 : 3 = 10$ cm.
}
\end{ex}

\begin{ex}
Một hình chữ nhật có chiều rộng 4 cm, chiều dài gấp 3 lần chiều rộng. Chiều dài hình chữ nhật đó là bao nhiêu xăng-ti-mét?
\shortans{12}
\loigiai{
$4 \times 3 = 12$ cm.
}
\end{ex}

\begin{ex}
Một hình chữ nhật có chiều rộng 4 cm, chiều dài 12 cm. Chu vi hình chữ nhật đó là bao nhiêu xăng-ti-mét?
\shortans{32}
\loigiai{
$(12 + 4) \times 2 = 32$ cm.
}
\end{ex}

\begin{ex}
Hình nào có 4 cạnh bằng nhau và 4 góc vuông? (Hình vuông ghi 1, hình chữ nhật ghi 2).
\shortans{1}
\loigiai{
Đó là hình vuông.
}
\end{ex}

\begin{ex}
Tổng số cạnh của 3 hình tam giác là bao nhiêu?
\shortans{9}
\loigiai{
$3 \times 3 = 9$ cạnh.
}
\end{ex}

\begin{ex}
Tổng số đỉnh của 4 hình tứ giác là bao nhiêu?
\shortans{16}
\loigiai{
$4 \times 4 = 16$ đỉnh.
}
\end{ex}

\begin{ex}
Một hình chữ nhật có chiều dài 15 cm, chiều rộng kém chiều dài 5 cm. Chiều rộng hình chữ nhật đó là bao nhiêu xăng-ti-mét?
\shortans{10}
\loigiai{
$15 - 5 = 10$ cm.
}
\end{ex}

\begin{ex}
Một miếng bìa hình vuông có chu vi 32 cm. Cạnh miếng bìa dài bao nhiêu xăng-ti-mét?
\shortans{8}
\loigiai{
$32 : 4 = 8$ cm.
}
\end{ex}

\begin{ex}
Một khu vườn hình vuông có chu vi là 40 m. Cạnh khu vườn đó dài bao nhiêu mét?
\shortans{10}
\loigiai{
$40 : 4 = 10$ m.
}
\end{ex}
